
\begin{tcolorbox}[
    enhanced,
    colback=white,
    colframe=black,
    arc=5mm,
    boxrule=1pt,
    left=2mm,
    right=2mm,
    top=1mm,
    bottom=1mm
]

\textbf{Task.} You are a data-quality auditor. Given a single reference passage and a single question, decide to what extent the passage alone contains the information needed to answer that question.

\textbf{Scoring rubric (integer 1--5).}\\
\textbf{5 — Fully answerable.} Every piece of information required by the question is explicit in (or straightforwardly derivable from) the passage; no outside knowledge needed.\\
\textbf{4 — Mostly answerable.} The passage supplies the core answer; only trivial world-knowledge or definitions lie outside.\\
\textbf{3 — Partially answerable.} The passage touches the topic but is missing a non-trivial component of the answer.\\
\textbf{2 — Minimally answerable.} Only distantly related evidence is present; a reliable answer cannot be drawn solely from this passage.\\
\textbf{1 — Unanswerable.} The passage provides no information that bears on the question.\\

\textbf{Output format.} First write two--three sentences of justification that quote the relevant span(s) of the passage. Then, on the final line, write exactly: \texttt{Rating: <integer 1-5>}.

\textbf{Demonstration.}\\
\emph{Passage:} Singapore is a sovereign city-state in maritime Southeast Asia... Modern Singapore was founded in 1819 by Sir Stamford Raffles as a British trading post.\\
\emph{Question:} In what year did Singapore gain independence?\\
\emph{Justification:} The passage mentions the 1819 founding by Raffles but says nothing about independence, which historically occurred later. No supporting span can be cited.\\
\emph{Rating: 1}

\textbf{Now audit the following real instance.}\\
\textbf{Passage:}\\
\{DOCUMENT\}\\

\textbf{Question:}\\
\{QUESTION\}\\

\end{tcolorbox}
